\documentclass[12pt,a4paper]{article}

\usepackage[a4paper,margin=2.5cm]{geometry}
\usepackage[T1]{fontenc}
\usepackage[utf8]{inputenc}
\usepackage[english]{babel}
\usepackage{setspace}
\usepackage{parskip}
\usepackage{titlesec}
\usepackage{hyperref}
\usepackage{enumitem}
\usepackage{amsmath,amssymb}
\usepackage{booktabs}
\usepackage{tabularx}
\usepackage{longtable}
\usepackage{array}
\usepackage{pgfgantt}
\usepackage{xcolor}
\usepackage{fancyhdr}

\hypersetup{colorlinks=true,linkcolor=black,urlcolor=blue,citecolor=blue,pdftitle={University Style Proposal EN}}
\setlength{\headheight}{15pt}
\pagestyle{fancy}
\fancyhf{}
\lhead{Research Proposal}
\rhead{Aging-Aware Adaptive Scheduling}
\cfoot{\thepage}

\definecolor{deepblue}{RGB}{0,51,102}
\definecolor{accentblue}{RGB}{0,102,204}
\titleformat{\section}{\large\bfseries\color{deepblue}}{\thesection.}{0.6em}{}
\titleformat{\subsection}{\normalsize\bfseries\color{accentblue}}{\thesubsection.}{0.6em}{}

\begin{document}

\begin{titlepage}
\begin{center}
{\Large \textbf{Izmir Institute of Technology}}\\[0.8em]
{\large Faculty of Engineering}\\
{\large Computer Engineering}\\[2.5em]

{\LARGE \textbf{THESIS / PROJECT PROPOSAL}}\\[1.0em]
{\Large \textbf{AI-Driven Aging-Aware Adaptive Time Quantum Scheduling}}\\[0.4em]
{\large Reliable and Efficient Scheduler Design with Cross-Layer Reliability Feedback}\\[2.0em]

\begin{tabular}{>{\bfseries}p{4.8cm}p{8.8cm}}
Student Name: & Omer Faruk Teker \\
Student ID: & To be specified \\
Supervisor: & To be specified \\
Program: & BSc \\
Department: & Computer Engineering \\
Submission Date: & April 2026 \\
\end{tabular}

\vfill
\begin{minipage}{0.92\textwidth}
\small
Note: Bracketed fields are left as placeholders for university, student, and supervisor information.
This version is prepared as a formal academic proposal with in-text citations and references.
\end{minipage}
\end{center}
\end{titlepage}

\pagenumbering{roman}
\tableofcontents
\newpage

\section*{Abstract}
\addcontentsline{toc}{section}{Abstract}

This research proposal aims to unify two lines of work that are usually treated separately:
\textit{adaptive time quantum selection} in CPU scheduling and \textit{aging-aware runtime control}
in hardware reliability management. Existing Round Robin variants mainly optimize waiting time,
turnaround time, and context-switch overhead through workload-aware quantum adaptation, while
reliability-aware systems focus on thermal management, DVFS, task remapping, and online monitoring
to preserve hardware lifetime \cite{behera2010,noon2011,marr2022,srinivasan2004,kumar2011}.
However, only limited work directly incorporates aging and thermal feedback into scheduler-level
time quantum control.

The proposed research introduces an \textit{aging-aware adaptive scheduler} that combines queue-state,
thermal telemetry, and aging estimates within a single control loop. The methodology is organized in
three stages: (i) heuristic aging-penalized quantum selection, (ii) physics-informed supervised learning
for quantum prediction, and (iii) contextual bandit or reinforcement learning for long-horizon policy
optimization. The core hypothesis is that such a scheduler can achieve either lower aging accumulation
at comparable performance or higher performance under the same aging budget.

\textbf{Keywords:} adaptive scheduling, Round Robin, time quantum, hardware aging,\\ reliability-aware control, reinforcement learning

\newpage
\pagenumbering{arabic}
\onehalfspacing

\section{Introduction}

Time quantum selection is one of the key design decisions in time-shared schedulers. When the quantum is
too small, context-switch overhead increases and throughput may drop. When it is too large, response time
and interactivity suffer. For this reason, a large body of dynamic Round Robin literature has proposed
workload-aware quantum selection strategies based on mean, median, quartile, or percentile statistics
\cite{behera2010,noon2011,marr2022,zohora2024}.

At the same time, scheduler decisions indirectly affect hardware lifetime. Imbalanced core usage, persistent
hotspots, long periods of high utilization, and strong thermal gradients can accelerate degradation mechanisms
such as NBTI, PBTI, HCI, electromigration, and thermal cycling \cite{srinivasan2004,agarwal2007,kumar2011}.
Despite this connection, scheduling literature and reliability-aware runtime control literature have largely
evolved independently.

This proposal addresses that gap by asking whether a scheduler can select time quantum values using not only
queue-state information but also thermal and aging feedback.

\section{Problem Statement}

The central problem of this research can be stated as follows:

\begin{quote}
Fixed or purely queue-statistics-driven dynamic time quantum selection does not explicitly account for the
long-term hardware aging and reliability cost of scheduling decisions. Therefore, a new scheduler-level control
mechanism is needed to jointly optimize performance, fairness, thermal behavior, and projected lifetime.
\end{quote}

This problem is particularly important in multicore systems, where task placement, context-switch rhythm, and
core utilization asymmetry have temporal and spatial effects on degradation behavior.

\section{Research Questions and Hypotheses}

\subsection{Research Questions}

\begin{enumerate}[leftmargin=*]
\item Can time quantum selection achieve a better performance-reliability trade-off when queue-state and thermal/aging telemetry are used together?
\item Can an aging-aware adaptive scheduler produce lower aging scores than fixed RR and classical dynamic RR under similar throughput levels?
\item Can physics-informed supervised learning or RL-based policies generalize better than heuristic policies under workload variation?
\item Does extending the scheduler action space with migration and DVFS lead to measurable projected lifetime improvements?
\end{enumerate}

\subsection{Hypotheses}

\begin{enumerate}[leftmargin=*]
\item Queue-aware and aging-aware hybrid schedulers provide a better latency-reliability trade-off than traditional dynamic RR schedulers.
\item Incorporating thermal and aging state into quantum selection reduces hotspot behavior and cumulative stress imbalance.
\item Physics-informed supervised models can select better time quanta than purely statistical schedulers.
\item RL-based control can reduce long-term aging cost when combined with suitable rewards and safety constraints.
\end{enumerate}

\section{Significance and Expected Contribution}

This research is significant for four reasons:

\begin{enumerate}[leftmargin=*]
\item It bridges scheduling literature and aging-aware reliability literature.
\item It treats time quantum selection as both a performance problem and a hardware health problem.
\item It compares heuristic, supervised learning, and RL-based methods in a shared simulation environment.
\item It directly tests the claims of ``same performance, less aging'' and ``same aging budget, better performance.''
\end{enumerate}

\section{Literature Review}

\subsection{Adaptive Time Quantum Scheduling}

Dynamic quantum selection has long been studied in the context of Round Robin scheduling. Behera et al.
reported improvements with a re-adjusted dynamic quantum policy \cite{behera2010}. Noon et al.
proposed a mean-based quantum computation strategy \cite{noon2011}, while more recent work has explored
median-average and percentile-based approaches \cite{marr2022,zohora2024}. Although these studies demonstrate
the value of workload-aware quantum selection, they generally omit thermal and aging effects.

\subsection{AI-Driven Scheduling}

Machine learning is increasingly used in scheduling research, but direct AI-driven CPU quantum selection remains
relatively sparse. Chitaliya et al.
framed quantum optimization as a prediction problem and evaluated ML-based quantum estimation \cite{chitaliya2023}.
Still, few works directly address scheduler-level reliability-aware quantum control.

\subsection{Hardware Aging and Reliability-Aware Control}

Srinivasan et al.
established the importance of lifetime reliability-aware microprocessor design \cite{srinivasan2004}. Agarwal et al.
highlighted the role of failure prediction in transistor aging management \cite{agarwal2007}, while Kumar et al.
developed adaptive techniques to compensate for aging-induced performance degradation \cite{kumar2011}. Reliability-aware
energy management and RL-based DVFS demonstrate that runtime performance-lifetime trade-offs can indeed be controlled
at higher system layers \cite{aydin2009,yeganeh2021}.

\subsection{Research Gap}

The key gaps in the literature are the following:

\begin{itemize}[leftmargin=*]
\item Quantum selection studies do not typically use hardware aging information.
\item Aging-aware runtime control studies rarely optimize scheduler quantum directly.
\item AI-based scheduler research rarely adopts a reliability-centric objective.
\end{itemize}

This proposal is positioned precisely at this intersection.

\section{Proposed Methodology}

\subsection{General Framework}

The proposed system consists of four components:

\begin{enumerate}[leftmargin=*]
\item Workload generator or trace source
\item Scheduler engine
\item Thermal estimator and aging estimator
\item Logging, evaluation, and analysis layer
\end{enumerate}

The policy layer will use queue-state, thermal telemetry, and aging estimates to generate quantum,
migration, and optionally DVFS decisions.

\subsection{Method 1: Heuristic Aging-Penalized Quantum}

The first method is an interpretable baseline:

\[
Q_t = f(S^{queue}_t) - \lambda_1 g(S^{thermal}_t) - \lambda_2 h(S^{aging}_t)
\]

Here, $f$ produces a base quantum using mean, median, or quartile statistics, while $g$ and $h$ provide
penalty terms derived from temperature, hotspot severity, thermal gradient, cumulative stress, and estimated
slack loss.

\subsection{Method 2: Physics-Informed Supervised Learning}

In the second stage, optimal time quantum selection is formulated as a prediction problem. For each workload
window, candidate quantum values are evaluated through a combined objective:

\[
J = w_1 \cdot \text{Latency} + w_2 \cdot \text{ContextSwitchCost} + w_3 \cdot \text{ThermalCost} + w_4 \cdot \text{AgingCost}
\]

The quantum that minimizes $J$ becomes the training label. Models such as Random Forest, XGBoost, and LightGBM
will then be used to predict quantum values. The approach is called physics-informed because both labeling and
feature design explicitly encode thermal and aging meaning.

\subsection{Method 3: Contextual Bandit and Reinforcement Learning}

In the third stage, scheduling is treated as a sequential decision problem. The state space includes queue length,
remaining burst statistics, per-core utilization, current temperature, temperature gradient, cumulative aging score,
previous quantum values, and migration history. A representative reward function is:

\[
R_t = \alpha \cdot \text{Throughput}_t - \beta \cdot \text{ResponseTime}_t - \gamma \cdot \text{ContextSwitches}_t - \delta \cdot \text{HotspotPenalty}_t - \eta \cdot \text{AgingIncrement}_t
\]

Algorithms such as DQN, Double DQN, or PPO may be used. All learned policies will be deployed under a safety filter
that enforces peak-temperature thresholds, starvation prevention, and migration-rate limits.

\section{Experimental Design}

\subsection{Simulation Environment}

The initial research platform will be a Python-based discrete-event scheduler simulator. This choice supports rapid
experimentation, reproducibility, and fair comparison of policy variants. Integration with gem5 or Sniper may be
considered in later phases.

\subsection{Workload Classes}

\begin{longtable}{@{}p{3.3cm}p{2.1cm}p{8.0cm}@{}}
\caption{Planned workload classes.}\\
\toprule
\textbf{Class} & \textbf{Type} & \textbf{Description} \\
\midrule
\endfirsthead
\multicolumn{3}{l}{\small\itshape (continued)}\\
\toprule
\textbf{Class} & \textbf{Type} & \textbf{Description} \\
\midrule
\endhead
\bottomrule
\endfoot

Short CPU-bound & Synthetic & Short bursts; high response-time sensitivity \\
Long CPU-bound & Synthetic & Long bursts; strong sensitivity to quantum size \\
Heavy-tail bursts & Synthetic & Pareto-like burst distribution \\
Bursty arrivals & Synthetic & Clustered arrivals that create queue pressure \\
Periodic tasks & Synthetic & Real-time-like periodic workloads \\
Hot-core stress & Scenario & Workload shaped to accumulate stress on one core \\
Balanced multicore mix & Scenario & Suitable for observing wear-leveling behavior \\
Trace-driven mix & Realistic & Derived from scheduler traces if available \\
\end{longtable}

\subsection{Baselines}

\begin{enumerate}[leftmargin=*]
\item Fixed Round Robin
\item Mean-based dynamic RR
\item Median/quartile-based dynamic RR
\item Thermal-aware scheduler
\item Aging-aware but non-AI scheduler
\item Proposed heuristic hybrid scheduler
\item Proposed supervised scheduler
\item Proposed bandit / RL scheduler
\end{enumerate}

\subsection{Thermal and Aging Models}

The thermal model is defined in a simple but consistent form:

\[
T_{t+1} = aT_t + bU_t + cS_t + dT_{amb}
\]

where $U_t$ is utilization, $S_t$ is a switching proxy, and $T_{amb}$ is ambient temperature.
The aging increment is modeled as:

\[
\Delta A_t = k_1 e^{\rho T_t} + k_2 U_t^p + k_3 C_t
\]

where $C_t$ captures additional stress related to context-switch or switching-rate behavior.
In later phases, BTI-like, HCI-like, EM-proxy, and thermal-cycling components may be tracked separately.

\section{Evaluation Metrics}

\subsection{Scheduler Performance}

\begin{itemize}[leftmargin=*]
\item Average waiting time
\item Average turnaround time
\item Average response time
\item Throughput
\item Context-switch count
\item Fairness and starvation indicators
\end{itemize}

\subsection{Reliability and Thermal Metrics}

\begin{itemize}[leftmargin=*]
\item Peak temperature
\item Mean temperature per core
\item Thermal gradient
\item Thermal-cycling severity
\item Cumulative aging score
\item Projected lifetime proxy
\end{itemize}

\subsection{Primary Comparative Claims}

The proposal directly tests two core claims:

\begin{itemize}[leftmargin=*]
\item lower aging accumulation at comparable performance
\item better performance under the same aging budget
\end{itemize}

\section{Illustrative Scenario and Objective Expectation Framework}

This section is added to make the proposal more concrete under a realistic operating scenario and to present both
the expected benefits and the possible limitations in an objective way.

\subsection{Illustrative Scenario}

Consider a four-core system with a mixed ready queue containing both short and long CPU-bound jobs. Assume that,
within a given scheduling window, the following conditions hold:

\begin{itemize}[leftmargin=*]
\item Core 1 has been under high utilization for a long period.
\item Core 1 temperature is higher than the others.
\item The aging score of Core 1 is increasing faster than the rest.
\item Queue occupancy is high, meaning that the scheduler is under performance pressure.
\end{itemize}

In this situation, a conventional dynamic RR scheduler would likely increase the quantum in order to reduce queue
pressure, but because it does not incorporate hardware health information, it may further intensify the stress on the
same core. By contrast, the proposed aging-aware scheduler would evaluate one or more of the following actions jointly:

\begin{enumerate}[leftmargin=*]
\item increase the quantum according to queue pressure, but limit that increase via thermal/aging penalties
\item migrate jobs to a less stressed core
\item use DVFS to move the stressed core into a more conservative operating mode if needed
\item filter migration or quantum decisions if fairness constraints are threatened
\end{enumerate}

This scenario makes clear that the proposed system is intended to answer not only ``which job should run and for how long,''
but also ``what is the hardware-health consequence of this scheduling decision?''

\subsection{Potential Positive Outcomes}

Under such a scenario, the following positive outcomes may be expected:

\begin{itemize}[leftmargin=*]
\item reduction in hotspot behavior on specific cores
\item more balanced core-wise aging score distribution
\item lower aging accumulation at comparable throughput relative to fixed RR or queue-only dynamic RR
\item lower thermal gradients and peak temperatures
\item improved projected lifetime proxy in long-running experiments
\end{itemize}

\subsection{Potential Negative Outcomes and Limitations}

For an objective evaluation, the proposal must also acknowledge the possibility of limited or negative outcomes:

\begin{itemize}[leftmargin=*]
\item migration and policy-computation overhead may increase total latency
\item overly aggressive aging protection may hurt throughput or turnaround time
\item thermal telemetry alone may not fully represent aging behavior, making the proxy model incomplete
\item RL-based policies may be sensitive to training data and reward weights
\item benefits observed in simulation may shrink in real-system deployment
\item for some workload families, classical dynamic RR may already perform well enough that additional gains are small
\end{itemize}

\subsection{Objective Success Criteria}

The work should be considered successful only if it can demonstrate outcomes such as:

\begin{enumerate}[leftmargin=*]
\item statistically meaningful aging-score reduction under comparable throughput
\item lower waiting time or turnaround time under the same aging budget
\item clear reduction in peak temperature or hotspot count
\item consistent improvement across multiple workload classes
\end{enumerate}

If these criteria are not fully satisfied, the study may still make a valuable contribution by clearly identifying the
conditions under which scheduler-level aging-aware control is beneficial and the conditions under which overhead dominates.

\section{Expected Challenges and Risk Mitigation}

\begin{table}[h!]
\centering
\renewcommand{\arraystretch}{1.3}
\begin{tabularx}{\textwidth}{lXX}
\toprule
\textbf{Risk} & \textbf{Issue} & \textbf{Mitigation Strategy} \\
\midrule
Simplified aging model & Limited physical realism & Sensitivity analysis, alternative proxies, parameter sweeps \\
Synthetic workload bias & Incomplete representation of real behavior & Trace-driven workloads and multiple distribution families \\
RL instability & Reward sensitivity and training difficulty & Establish heuristic and supervised baselines first \\
Policy overhead & Control cost may offset benefits & Discrete action space and sparse update intervals \\
Fairness degradation & Aging protection may hurt some tasks & Fairness penalties and starvation bounds \\
\bottomrule
\end{tabularx}
\end{table}

\section{Work Plan and Timeline}

\begin{center}
\begin{ganttchart}[
    hgrid,
    vgrid,
    time slot format=isodate-yearmonth,
    x unit=0.95cm,
    y unit title=0.6cm,
    y unit chart=0.7cm,
    title/.append style={fill=deepblue!10},
    bar/.append style={fill=accentblue!60},
    bar height=0.55
]{2026-05}{2026-11}
\gantttitlecalendar{year, month}\\
\ganttbar{Literature consolidation and proposal finalization}{2026-05}{2026-05}\\
\ganttbar{Baseline scheduler simulator}{2026-05}{2026-06}\\
\ganttbar{Thermal and aging model integration}{2026-06}{2026-07}\\
\ganttbar{Heuristic hybrid scheduler}{2026-07}{2026-08}\\
\ganttbar{Dataset generation and labeling}{2026-08}{2026-09}\\
\ganttbar{Supervised learning experiments}{2026-09}{2026-10}\\
\ganttbar{Bandit / RL experiments}{2026-10}{2026-11}\\
\ganttbar{Ablation, analysis, and final writing}{2026-11}{2026-11}
\end{ganttchart}
\end{center}

\section{Expected Outcomes}

The expected outcome is that the proposed scheduler will demonstrate a better performance-reliability trade-off
than conventional schedulers. The strongest experimental outcomes would be:

\begin{enumerate}[leftmargin=*]
\item significant aging-score reduction under comparable throughput
\item lower waiting/turnaround time under the same aging limit
\item improved peak temperature and hotspot behavior
\item more consistent policy behavior across workload families
\end{enumerate}

However, an objective evaluation must also report the following possibilities:

\begin{itemize}[leftmargin=*]
\item gains may be small for some workloads
\item some policies may reduce aging while worsening latency
\item the best-performing policy and the most reliability-friendly policy may not be the same
\item the RL-based version may not always outperform heuristic or supervised variants
\end{itemize}

Therefore, the purpose of the study is not merely to claim that the proposed method is universally best, but rather to
show under which conditions it provides advantages, under which conditions it introduces cost, and how the trade-off
surface is shaped.

\section{Conclusion}

This proposal presents a detailed and feasible research framework for bridging adaptive time quantum scheduling and
aging-aware runtime control. The proposed system does not treat quantum selection purely as a performance problem;
it explicitly incorporates hardware health into the decision process. In this sense, the work has the potential to
make a meaningful contribution at the intersection of operating systems, computer architecture, and AI for systems.

\newpage
\begin{thebibliography}{99}

\bibitem{behera2010}
H. S. Behera, R. K. Mohanty, and D. Nayak,
``A New Proposed Dynamic Quantum with Re-Adjusted Round Robin Scheduling Algorithm and Its Performance Analysis,''
\textit{International Journal of Computer Applications}, 2010.

\bibitem{noon2011}
A. Noon, A. Kalakech, and S. Kadry,
``A New Round Robin Based Scheduling Algorithm for Operating Systems: Dynamic Quantum Using the Mean Average,''
arXiv preprint arXiv:1111.5348, 2011.

\bibitem{marr2022}
Sakshi, C. Sharma, S. Sharma, S. Kautish, S. A. M. Alsallami, E. M. Khalil, and A. W. Mohamed,
``A New Median-Average Round Robin Scheduling Algorithm,''
\textit{Alexandria Engineering Journal}, 2022.

\bibitem{zohora2024}
M. F. Zohora, F. Farhin, and M. S. Kaiser,
``An Enhanced Round Robin Using Dynamic Time Quantum for Real-Time Asymmetric Burst Length Processes,''
\textit{PLOS ONE}, 2024.

\bibitem{chitaliya2023}
P. Chitaliya et al.,
``Time Quantum Optimization in Round Robin Algorithm,''
\textit{NMITCON}, 2023.

\bibitem{srinivasan2004}
J. Srinivasan et al.,
``The Case for Lifetime Reliability-Aware Microprocessors,''
\textit{ISCA}, 2004.

\bibitem{agarwal2007}
M. Agarwal et al.,
``Circuit Failure Prediction and Its Application to Transistor Aging,''
\textit{VTS}, 2007.

\bibitem{aydin2009}
H. Aydin and D. Zhu,
``Reliability-Aware Energy Management for Periodic Real-Time Tasks,''
\textit{IEEE Transactions on Computers}, 2009.

\bibitem{kumar2011}
S. V. Kumar, C. H. Kim, and S. S. Sapatnekar,
``Adaptive Techniques for Overcoming Performance Degradation Due to Aging in CMOS Circuits,''
\textit{IEEE Transactions on VLSI Systems}, 2011.

\bibitem{yeganeh2021}
A. Yeganeh-Khaksar et al.,
``Ring-DVFS: Reliability-Aware Reinforcement Learning-Based DVFS for Real-Time Embedded Systems,''
\textit{IEEE Embedded Systems Letters}, 2021.

\end{thebibliography}

\end{document}
